% !TeX root = ../main.tex

\chapter{基于融合特征的分类方法}

上一章介绍了基于BYOL模型的流级特征分类方法，提取了流级别的特征表示。然而，仅依靠流级别特征难以全面捕捉网络流量的复杂模式，如图~\ref{fig:90f}所示：在仅提取16个数据包长度序列的情况下，VNAT数据集中的C2和file类别存在混淆情况，原因是两种流量在流级别的特征表示非常相似。为了进一步提升分类性能，本章将介绍基于n-gram技术的包级别特征提取方法，并设计了基于注意力机制的特征融合方法，将流级别和包级别的特征进行融合，构建一个更为有效的加密流量分类模型。
\begin{figure}[ht]
  \centering
  \includegraphics[width=1\textwidth]{90f.png}
  \bicaption{混淆情况}{}
  \label{fig:90f}
\end{figure}

\section{系统模型}
本章提出了一种基于双模态融合特征的深度学习网络流量分类模型。该模型旨在解决单一维度特征在面对复杂加密流量时表达能力不足的问题。系统整体架构由三个核心阶段组成。首先是包级别特征的增强型 N-gram 序列建模，用于挖掘报文内部的字节结构信息和文本顺序信息；其次是基于交叉注意力机制的特征融合模块，将流级别特征与包级别特征进行非线性对齐与融合，利用交叉注意力机制挖掘两种维度特征间的相关性，并在训练过程中动态地调节两种特征的权重从而达到更好的融合效果。最后，我们还设计了深度残差分类器MLP分类器用于进行分类任务，分类器主体采用了由多个残差块堆叠形成的网络结构，每个残差块内部均包含全连接层、批归一化与 Dropout 结构，通过引入跳跃连接机制，有效缓解了深层网络中的梯度消失问题，并提升了模型的表达能力和泛化性能，确保模型能够精准捕获复杂的流量分类边界并最终输出准确的流量类别。
\section{包级别特征提取}
在传统的基于数据包内容的分类方法中，通常依赖数据包的有效载荷作为核心识别特征。然而，有效载荷中往往包含用户身份、通信内容及行为偏好等高度敏感的信息，直接将其作为深度学习模型的输入，极易引发潜在的隐私泄露风险与伦理争议。此外，随着各种加密协议的普及，有效载荷呈现高度熵增的伪随机状态，导致传统的载荷分析手段较难达到准确的分类效果。

Hu 等人的研究发现，虽然加密协议屏蔽了应用层内容，但网络层与传输层的数据包首部数据仍保留了丰富的协议交互逻辑与特性。其实验结果表明，仅通过提取报文头部的字节序列进行建模，即可在多种复杂业务场景下实现与全载荷方案相媲美的分类精度。基于此，本研究在确保用户隐私边界的前提下，将有效载荷数据提出，仅使用 IP 数据报头部数据作为模型的原始输入。通过对报文头部进行结构化特征挖掘，旨在构建一种既符合隐私保护要求，又能适配加密网络环境的分类模型。

\subsection{报文首部截取与预处理}
为了捕获网络协议交互的核心语义并构建特征输入，我们首先对原始流量进行流分割，基于五元组信息{源 IP、目的 IP、源端口、目的端口及传输层协议}将捕获的数据包切分为逻辑连接的数据流。接着预先截取每个流前5个数据包的 IP 报文首部作为原始特征输入。

在特征空间构建过程中，需针对报文首部的结构特性进行处理。标准的IPV4首部包含20字节的固定长度字段，其结构如表~\ref{tab:IPheader}所示。IP分组报头涵盖了段版本、首部长度、优先级与服务类型、总长度、标识、标志、段偏移、生存时间、协议、校验和、源IP地址和目的IP地址这12个关键字段。然而，源/目的IP地址字段与特定的网络拓扑配置高度耦合，其数值分布更多反映了网络部署环境而非业务流本身的协议特征。从网络流量分类任务的角度来看，这两个字段属于冗余信息，会影响模型的泛化能力和分类结果。因此，本工作在预处理阶段将源/目的IP地址字段视为剔除，从而实现特征降维与关键语义的聚焦。

\begin{table}[h]
    \centering
    \bicaption{IP数据报头结构}{}
    \label{tab:IPheader}
    \begin{tabular}{|c|c|c|c|c|}
        \hline
        \textbf{0-3 bit} & \textbf{4-7 bit} & \textbf{8-15 bit} & \textbf{16-18 bit} & \textbf{19-31 bit} \\ \hline
        版本 (Version) & 首部长度 (IHL) & 服务类型 (ToS) & \multicolumn{2}{c|}{总长度 (Total Length)} \\ \hline
        \multicolumn{3}{|c|}{标识 (Identification)} & 标志 (Flags) & 片偏移 (Fragment Offset) \\ \hline
        生存时间 (TTL) & 协议 (Protocol) & \multicolumn{3}{c|}{首部校验和 (Header Checksum)} \\ \hline
        \multicolumn{5}{|c|}{源 IP 地址 (Source IP Address)} \\ \hline
        \multicolumn{5}{|c|}{目的 IP 地址 (Destination IP Address)} \\ \hline
        \multicolumn{4}{|c|}{可选字段 (Options，若有)} & 填充 (Padding) \\ \hline
    \end{tabular}
\end{table}

\subsection{N-gram嵌入处理}
在现有的基于数据包内容的分类方法中，主流方法通常将报文视为独立的字节流，并将每个字节视为离散的输入特征。然而，IP 数据包的字节序列具有比自然语言更为严谨的结构化语义。报文首部中的多个连续字节通常共同构成一个具有明确物理含义的字段，例如总长度字段由两个字节组成，段偏移字段涉及位运算关联。此外，报文首部字节间存在固定的位置关系，这些结构化信息与顺序特征对于识别协议模式至关重要。融入这些信息可以增强特征表示。为了充分挖掘字节间的结构化信息与顺序特征，本研究引入了N-gram嵌入方案。通过在预处理后的报文首部滑动取窗，提取具备上下文关联的语义特征。

本工作使用大小为 $n \in \{1, 2, 3\}$ 的滑动窗口对报头数据进行扫描。将 1 字节、2 字节和 3 字节的段作为特征生成的标记，所有窗口的步长均为 1。对于预处理过程中获得的IP数据包报头，1-gram 捕捉孤立的字节信息，而 2-gram 和 3-gram 通过构建长度分别为 2 和 3 的字节组合，提取报文内部的结构化信息与顺序特征。最后将得到的嵌入式特征按顺序连接起来，形成最终的包级特征。整个过程如图 3 所示。

\begin{figure}[ht]
  \centering
  \includegraphics[width=1\textwidth]{ngram.png}
  \bicaption{n-gram嵌入处理}{}
  \label{fig:ngram}
\end{figure}

对于1-gram得到的字节形式特征，我们对其字节流进行归一化处理。以此消除量纲差异并加速深度学习模型的梯度收敛。本工作将原始字节在 $[0, 255]$ 范围内的取值线性映射至 $[-1, 1]$ 的数值区间。这种处理方式在保持报文首部关键字段原始分布特征的同时，通过中心化与缩放操作提升了线性特征的分布稳定性，为后续通过非线性模型提取包级别特征奠定基础。

而对于2-gram和3-gram得到的字节组合特征，我们将其十六进制字节组合转换为对应的十进制数值 $ngram_{dec}$。通过这种映射，原始的离散字节被转换至具备数值统计意义的高维特征空间中。由于网络协议中存在诸如标识符、校验和或高精度时间戳等字段，其产生的 $ngram_{dec}$ 往往呈现极高的随机性以及长尾分布。若采用传统的线性归一化方法，由于 n-gram 理论最大值 $16^{2n}-1$ 极大，会导致大部分关键信息被压缩至极小的数值区间内，产生严重的特征淹没效应。因此，我们引入了非线性对数变换来缓解这一问题。
$$
log\_val = \ln(1 + ngram_{dec})
$$
通过对 $ngram_{dec}$ 进行对数变换，模型能够在保持原始数值单调性的前提下，有效地压缩高值区间、扩张低值区间，同时保留特征之间的相对关系，从而缓解数据分布极度不均的问题。随后，利用理论最大值的对数映射 $log\_max = \ln(1 + max\_val)$ 对数据进行重缩放，将其映射至与1-gram特征相同的 $[-1, 1]$ 区间内，确保不同 n-gram 特征在同一数值范围内进行融合与学习。
$$
normalized = 2 \times \frac{log\_val}{log\_max} - 1
$$
这种处理方式不仅保留了数据报头之间的词组信息和各关键字之间的顺序信息，还有效地压缩了极端值的影响，使得模型能够更好地学习到数据包级特征中的细微差异，从而提升分类性能。

为了使后续下游分类任务能够正常进行，我们还将处理后的特征进行进行维度对齐，从而保证分类模型输入张量的维度一致性。针对通信持续时间较短、数据包总数不足 5 个的短流，模型通过计算现有长度来动态生成占位包。填充数值统一设定为 -1，以区分于正常数据包的特征值。这种填充方式不仅保证了多级特征按顺序连接后的特征向量长度一致，还通过显式的边界填充符号降低了填充数据对模型真实特征学习的干扰，帮助模型在训练过程中识别并适应不同长度流量的特征分布差异，从而提升模型的泛化能力和分类性能。

\section{交叉注意力特征融合机制}
在完成上述包级别特征与上一章流级别特征的提取后，我们还需要将两个模态的两类特征进行高效地融合，从而提升模型的分类性能。传统的特征融合方法多采用简单的向量拼接或加权求和，这种线性处理方式往往忽略了流级别特征与包级别特征之间深层的非线性关联，难以充分发掘不同模态特征在刻画特定网络行为时的互补潜力。为了深度挖掘这种关联信息，本研究构造了基于交叉注意力机制的特征融合模块。该模块通过模拟人类视觉注意力的分配逻辑，以流级别特征作为引导信号，动态地在包级别特征序列中检索并聚焦于关键的报文模式。通过在隐空间内进行特征重构与空间对齐，该机制不仅能够自动学习不同模态特征的贡献权重，还能有效缓解异构数据间的维度冲突，从而生成具备高度判别力的融合特征，为后续分类网络提供更加有助于流量分类的输入。

\subsection{线性投影与空间对齐}
由于流级特征 $F$ 与包级特征 $P$ 来源于不同的提取手段，其特征空间并不一致。模型首先通过投影矩阵 $W_f$ 和 $W_p$ 将其映射至 256 维的对齐空间：$$f_{proj} = \sigma(W_f \cdot F + b_f), \quad p_{proj} = \sigma(W_p \cdot P + b_p)$$其中 $\sigma$ 为线性对齐算子。

\subsection{交叉注意力流}
模型利用流级特征作为 Query ($Q$)，去检索包级特征中的 Key ($K$) 和 Value ($V$)：注意力权重分配：通过 8 个独立的注意力头计算包级特征对当前分类任务的重要性分布。残差融合：将注意力输出与原始流级投影特征进行残差连接，并经过层归一化（LayerNorm）处理：$$fused_{initial} = \text{LayerNorm}(\text{Attn}(Q, K, V) + f_{proj})$$前馈网络增强：通过一个隐藏层维度为 512 的双层 FFN 进一步挖掘跨模态特征的非线性组合。

\section{深度残差分类器设计 (Residual MLP)}
分类模块通过增加网络深度来提升模型对复杂流量模式的建模能力。

\subsection{残差块 (Residual Block) 内部构造}
每个残差块包含两个全连接层，并在每个线性层后嵌入 Batch Normalization (BN)：BN 层作用：加速训练过程中的收敛，缓解深层网络中的内部协变量偏移问题。Dropout 层：在每个块内部设置 0.2 的随机丢弃率，防止模型过度拟合训练集中的特定噪声。跳跃连接：块输入直接与第二个 BN 层的输出相加，使模型即使在 5 层残差结构（共计 10 层以上全连接）下仍能保持高效的梯度流动。

\subsection{输出层逻辑}
模型末端接一个全连接层，映射至 6 维输出空间，分别对应 $categories = [\text{socialapp, chat, email, file, streaming, VoIP}]$。

\section{训练方法}

\subsection{学习率调度：OneCycleLR}
在融合模型的预训练阶段，采用了 OneCycleLR 调度策略：阶段性升降：学习率在前 30\% 的 Epoch 内从 $1\times 10^{-5}$ 升至 $2\times 10^{-3}$，随后按余弦函数逐渐衰减至 $1\times 10^{-8}$。优势：这种“超级收敛”策略有助于模型跳出局部最优解，在 10,000 个 Epoch 的训练中寻找全局最优点。

\subsection{验证驱动的回溯降温策略}
在主分类器训练期间，模型引入了更为严格的质量控制：
抖动检测：设置 $\text{decay\_patience} = 200$。若连续 200 代验证集准确率未提升，判定为陷入局部抖动。
回溯机制：自动加载历史最优模型参数（Best Model），并将当前学习率强制减半（乘以 0.5），实现“降温寻优”。
最终停机：若回溯多次后仍连续 1000 个 Epoch 无进展，则触发早停机制。
\subsection{评估指标设置}
模型通过计算 归一化混淆矩阵 (Normalized Confusion Matrix) 进行评估，通过对混淆矩阵按行求和并计算各类别预测比例，直观展示模型在易混淆类别（如 Chat 与 SocialApp）上的辨析能力。

\section{实验结果与分析}
本工作在4个公开数据集上对分类模型进行了评估，通过与FS-NET模型，1d-cnn模型两个单模态分类方法，以及双模态分类方法DM-HNN模型进行对比实验。同时，还将完整的模型分别与仅使用流级别特征和仅使用包级别特征的模型进行了消融实验对比。实验结果表明，我们提出的分类模型在多个数据集上均取得了优异的性能，显著提升了分类效果。同时，消融实验结果表明，流级别特征和包级别特征的融合能够显著提升分类性能，验证了两类特征的互补性。

\subsection{参数设置}
\subsection{对比算法}
\subsection{算法结果分析}

\section{本章小结}


